\section{The actual NStage blocks and a relative-adjoint reduction}
\label{sec:relative-nstage-blocks}

The exact theta transport of AC-G42 preserves a packet and its
transported lattice.  The coefficient needed for the cubic problem
is instead relative to two packets.  The following statements locate
that distinction inside the supplied source and give the integral
coefficient comparison that a new relative theorem would have to use.
They do not reverse any congruence-index inequality.

\begin{proposition}[AC-G44: literal blocks, local conditions, and cancellation]
\label{AC-G44}
Fix a source datum at a coefficient prime $\ell$, and write
$\mathcal O$ for its complete coefficient DVR, $L$ for its fraction
field, $D=L/\mathcal O$, $F$ for its host field, and
$\varepsilon$ for the $\ell$-adic cyclotomic character.
The source semisimple decomposition at its selected point is
\[
 r_\pi\simeq\rho_\pi\oplus\bigoplus_{c=1}^{t}a_c,
 \qquad a_c=a_1\varepsilon^{-(c-1)}.                 \tag{G44.1}
\]
Its good-block and scalar-ratio coefficients are respectively
\[
 W=\rho_\pi a_1^{-1},\qquad
 W^\vee\varepsilon,\qquad a_ca_1^{-1}
                   =\varepsilon^{-(c-1)}.          \tag{G44.2}
\]
More precisely, the source quantity \texttt{sPi} is the length
of the Selmer group of $D\otimes_{\mathcal O}W$ having its stated
Bloch--Kato local condition at $\ell$, no condition at $S\cup Q$,
and the unramified condition elsewhere.  Its separately supplied
\texttt{shaLen} is the strict Tate group: unramified outside
$\ell$ and locally zero at every prime above $\ell$.
These are different definitions.

Simultaneously replacing $r_\pi,\rho_\pi,a_c$ by
$\nu r_\pi,\nu\rho_\pi,\nu a_c$, for any integral unit-valued
character $\nu$, leaves every coefficient in (G44.2) and
$r_\pi a_1^{-1}$ exactly unchanged.  The corresponding local
cohomology groups and specified local subgroups are unchanged
when the same period data are used.

Even if the displayed formula for $a_c$ in (G44.1) were removed,
the additional source hypotheses
\texttt{SemisimpleDecomposition}, \texttt{admG4}, and
\texttt{TameGeneratorConditions} force every $a_c/a_1$ to be
unramified at every finite prime not above $\ell$.
Consequently, over any host unramified at $269$ and containing
$\mathbb Q(\omega)$, the coefficient
$\chi_{269}\varepsilon^{-e}$ cannot replace such a scalar ratio
at $\ell\ne269$ while retaining those hypotheses.
No assertion about the coefficient prime $269$ is made by this
inertia argument.

Finally, assume the residual good-block spanning condition in
\texttt{admG3}.  For every integral rank-one character $\eta$,
\[
 \operatorname{Hom}_{\mathcal O[G_F]}
       (D(\eta),D\otimes W)=0,
 \qquad
 \operatorname{Hom}_{\mathcal O[G_F]}
       (D\otimes W,D(\eta))=0.                       \tag{G44.3}
\]
Thus a coefficient inclusion or quotient of the actual good
block does not supply the scalar cubic group.  This statement
does not assert that arbitrary abstract maps between the resulting
cohomology groups must vanish.
\end{proposition}
\begin{proof}
\proofstep{1. Reading the precise coefficient and local definitions.}
All source locators in this section use the previously pinned
commit \texttt{7cbf7a0c32fca1badf05931879a37c3f8a2b03b4}.
In \path{Formalism/SharpRestriction/LevelLetterCocycles/DataN.lean},
\texttt{NData.chi\_def}, lines 126--128, is exactly the formula
for $a_c$ in (G44.1).  The same directory's
\path{DecompositionData.lean},
\texttt{SemisimpleDecomposition}, lines 267--271, is the
representation decomposition in that equation.
In \path{QuotientHomNonvanishingBound.lean},
\texttt{Glob.Wp} and \texttt{Glob.T2}, lines 45--51,
and their matrix versions \texttt{WpO}, \texttt{T2O},
\texttt{rTwO}, \texttt{rDuO}, \texttt{chiRatO},
lines 271--284, give precisely (G44.2).

The same file defines \texttt{selG} in lines 786--789 and
\texttt{Br.selPi}, \texttt{Br.sPi}, \texttt{Br.sFin} in
lines 857--864.  Its \texttt{LlG}, lines 779--784, supplies
the local Bloch--Kato condition at $\ell$.
The complete case distinction is
\path{Formalism/GaloisCohomology/Statement.lean},
\texttt{SelAll} and \texttt{bkCond}, lines 158--168:
at $\ell$ use the supplied local submodule, away from $\ell$
use the whole local $H^1$ on $S\cup Q$, and otherwise use
\texttt{unrAt}.
In that file, \texttt{LData.hS}, lines 52--64,
\texttt{qzTw}, lines 67--69, and \texttt{Sha},
\texttt{shaLen}, lines 75--78, instead define the strict
Tate group on the Galois group unramified outside $\ell$.
The length identification with the row datum is the separate
field \texttt{NData.sha\_eq}, lines 66--68 of \texttt{DataN.lean}.
In particular, the notation \texttt{sPi} does not identify the
source good-block Selmer group with that strict group.

Cancellation under simultaneous twisting is an equality of
integral matrices:
$(\nu\rho)(\nu a_1)^{-1}=\rho a_1^{-1}$ and
$(\nu a_c)(\nu a_1)^{-1}=a_ca_1^{-1}$.
Taking a dual and then tensoring with $\varepsilon$ preserves
the equality.  Thus this calculation does not require a
rational projector, a change of lattice, or a claim about
invariance of Bloch--Kato conditions under an arbitrary twist:
the representations to which those conditions are applied are
literally equal.

\proofstep{2. The independent inertia obstruction.}
Let $v\nmid\ell$.  If $v\in S\cup Q$, the second clause of
\texttt{TameGeneratorConditions},
\texttt{DecompositionData.lean}, lines 224--230, says
$a_c(g)=a_1(g)$ for every $g\in I_v$ and every $1\le c\le t$.
Their ratio is therefore trivial on $I_v$.
If $v\notin S\cup Q$, \texttt{admG4}, lines 198--201,
says that $r_\pi(g)=1$ for $g\in I_v$.
The direct-sum decomposition forces every scalar constituent,
including $a_c$ and $a_1$, to be one on that inertia group.
This proves the assertion at every $v\nmid\ell$.

For the stated cubic host, the valuation of $269$ at every prime
above $269$ is one.  The polynomial $X^3-269$ is Eisenstein at
that prime.  Since the host contains $\mu_3$ and the residue
characteristic is not three, the local Kummer extension is
cyclic, totally tamely ramified, and of degree three.
Consequently $\chi_{269}|_{I_v}$ is nontrivial.
For $\ell\ne269$, the cyclotomic character is unramified at $v$:
roots of unity of order prime to the residue characteristic
belong to the maximal unramified extension, by the simple-root
lifting of their residue roots.  Thus
$\chi_{269}\varepsilon^{-e}|_{I_v}$ is still nontrivial.
This contradicts the preceding consequence of the source
hypotheses, independently of \texttt{chi\_def}.

\proofstep{3. The actual good rank and all torsion coefficient maps.}
The source \path{Formalism/LocalRingCardinality/Params.lean}
defines $r=k+m-1$, with even $k,m$, $m\ge2$, and
$k\ge m+3$.  Hence $k\ge6$ and $r\ge7$.
The final clause of \texttt{admG3},
\texttt{DecompositionData.lean}, lines 191--196, says that
the residual matrices of $\rho_\pi$ span the full matrix
algebra.  After any extension of the residue field they still
span the full matrix algebra, so their representation is
absolutely irreducible.  Multiplying each matrix by the
nonzero scalar $\overline a_1(g)^{-1}$ does not change its
linear span.  The reduction $\overline W$ is therefore also
absolutely irreducible, of dimension $r>1$.

Let $\varpi$ be a uniformizer and $k_{\mathcal O}$ the residue
field.  For each $n\ge1$, the finite module
$(D\otimes W)[\varpi^n]$ has a filtration whose successive
nonzero factors are $\overline W$: identify it with
$W/\varpi^n W$ by multiplication by $\varpi^{-n}$ and use
its $\varpi$-adic filtration.  In contrast,
$D(\eta)[\varpi^n]$ has only the one-dimensional factor
$k_{\mathcal O}(\overline\eta)$.
A nonzero finite image of an equivariant map in either
direction in (G44.3) would be a quotient of one of these
finite modules and a submodule of the other.  Jordan--H\"older
would give it a common simple factor.  There is none, since
one simple factor has dimension $r>1$ and the other has
dimension one.  Every element of either divisible module
is killed by some $\varpi^n$.  Applying this argument to
each restriction proves both vanishing statements.
\end{proof}

\begin{proposition}[AC-G44-trace: the integral mixed-adjoint reduction]
\label{AC-G44-trace}
Let $\Gamma$ be profinite, let $T$ be a continuous free
$\mathcal O$-representation of positive rank $r$, and let
$\eta:\Gamma\to\mathcal O^\times$ be continuous.  Put
\[
 A=D(\eta),\qquad
 B=D\otimes_{\mathcal O}\operatorname{Hom}_{\mathcal O}
                      (T,T(\eta)).                    \tag{G44.4}
\]
Fix any family of closed subgroups $\Gamma_v$ and define
$S(A),S(B)$ as the kernels of restriction to their local
$H^1$ groups.  Then there are natural maps
\[
 S(A)\xrightarrow{\ i\ }S(B)\xrightarrow{\ \mathrm{tr}\ }S(A),
 \qquad \mathrm{tr}\,i=r.                              \tag{G44.5}
\]
If $\ell\nmid r$, the first map is split injective.
In general, whenever the two quantities on the right are finite,
\[
 \operatorname{length}_{\mathcal O}S(A)
 \le \operatorname{length}_{\mathcal O}S(B)
       +\operatorname{length}_{\mathcal O}
                         H^1(\Gamma,A[r]).             \tag{G44.6}
\]
In particular that finiteness implies the finiteness of $S(A)$.
The same comparison holds for any additive normalization of
length, including the source's coefficient-normalized length.
The assertion also applies to functorial Selmer local conditions
for which both displayed coefficient maps preserve the chosen
local subgroups.
\end{proposition}
\begin{proof}
The action on an endomorphism is
$g\cdot X=\eta(g)T(g)XT(g)^{-1}$.
The scalar map $a\mapsto a\operatorname{Id}_T$ is equivariant,
as is $X\mapsto\operatorname{tr}(X)$, since trace is invariant
under conjugation.  Their composition is multiplication by $r$.
These maps are integral; no averaging over $\mu_3$ occurs.
They induce the maps on continuous cocycles, commute with
coboundaries and restriction to every $\Gamma_v$, and therefore
induce (G44.5) on the indicated kernels.
If $r$ is a unit, $r^{-1}\operatorname{tr}$ is a retraction.

In general $\ker i\subset S(A)[r]$ by (G44.5).
The exact sequence of discrete continuous $\Gamma$-modules
\[
                  0\longrightarrow A[r]\longrightarrow A
                   \xrightarrow{\ r\ }A\longrightarrow0
                                                               \tag{G44.7}
\]
is exact because $D=L/\mathcal O$ is divisible.
Its cohomology sequence shows that $H^1(\Gamma,A)[r]$
is an image of $H^1(\Gamma,A[r])$.
Thus $\ker i$ has length at most the second term in (G44.6),
while $\operatorname{im}i\subset S(B)$ has length at most the
first.  The exact sequence with this kernel and image gives
the claimed finiteness and inequality.  This argument uses
only finite-length additivity and therefore also proves the
normalized version.

For $\eta=\chi_{269}\varepsilon^{-e}$, the underlying rational
mixed representation is
\[
 \operatorname{Hom}(V,\chi_{269}V)(-e)
                 =\chi_{269}\varepsilon^{-e}\otimes\operatorname{End}(V).
                                                               \tag{G44.8}
\]
This identifies the coefficient to which the trace reduction
applies.  For number fields one may take $\Gamma=G_{F,\Sigma}$
with $\Sigma$ containing all coefficient ramification, and
include decomposition groups at $\ell$ and the desired inertia
subgroups among the $\Gamma_v$.  One must actually match those
conditions to a proposed relative theorem.  Formula (G44.6)
is not a claim that the source's \texttt{selPi} is $S(B)$,
and its finite correction is estimated arithmetically in
AC-G44-budget below.
\end{proof}

\begin{proposition}[AC-G44-center: what a balanced branch comparison sees]
\label{AC-G44-center}
Let $T_0\subset\mathcal O^X$ be the finite local algebra
of AC-G39 and fix $x\in X$.  Suppose its paired algebra is
\[
 C=T_0\times_{T_0/\delta T_0}T_0,
 \qquad \delta=\omega-1,                               \tag{G44.9}
\]
with the understanding that $\delta$ is a unit and the algebra
is the full product at $\ell\ne3$.
Let $\epsilon_0(a,b)=a(x)$ and $\epsilon_1(a,b)=b(x)$.
Give $\mathcal O$ and $D$ the corresponding $C$-actions.
Then evaluation at one gives an exact identification
\[
 \operatorname{Hom}_{C}(\mathcal O_{\epsilon_0},D_{\epsilon_1})
                     =D[\delta].                      \tag{G44.10}
\]
More generally, every $C$-linear map from a module on which
$C$ acts through $\epsilon_0$ to one on which it acts through
$\epsilon_1$ has image killed by $\delta$.
Thus these balanced maps vanish away from three.
For $\mathcal O=\mathbb Z_3[\omega]$, their images are killed
by the uniformizer $\lambda=\omega-1$.
Any such map out of the fixed-field strict cubic group of
AC-K23 has kernel containing $\lambda S_e$ and hence has
kernel length at least $m_e-1\in\{3,4\}$.
\end{proposition}
\begin{proof}
The ideal generated by all differences
$\epsilon_1(c)-\epsilon_0(c)$, $c\in C$, is exactly
$\delta\mathcal O$.  Every difference belongs to this ideal
by (G44.9), and $(0,\delta\cdot1)\in C$ has difference
$\delta$.  A $C$-linear map from $\mathcal O_{\epsilon_0}$
is determined by $m=f(1)\in D$ and satisfies
\[
       (\epsilon_1(c)-\epsilon_0(c))m=0\quad(c\in C).
\]
These conditions are equivalent to $\delta m=0$, and every
such $m$ defines a $C$-linear map.  This proves (G44.10).
For a general map $f$ and vector $v$, central linearity gives
$(\epsilon_1(c)-\epsilon_0(c))f(v)=0$; applying the specific
element $(0,\delta)$ proves the image assertion.
At coefficient three in the specified ring, AC-G22 and
AC-K23 give $S_e\simeq\mathcal O/\lambda^{m_e}$ with
$m_e=4+1_{e\equiv2\bmod3}$.  An image killed by $\lambda$
forces $\lambda S_e$ into the kernel.  This submodule has
length $m_e-1$, proving the final claim.
\end{proof}

\begin{corollary}[AC-G44-budget: the trace correction has small arithmetic support]
\label{AC-G44-budget}
Let $F$ be totally imaginary and contain $K=\mathbb Q(\omega)$.
For each prime $\ell$ let $\mathcal O_\ell$ be a finite coefficient
DVR containing the values of $\chi_{269}$, and put
$b_\ell=[\mathcal O_\ell:\mathbb Z_\ell]$ and
$A_{\ell,e}=(\operatorname{Frac}(\mathcal O_\ell)/
\mathcal O_\ell)(\chi_{269}\varepsilon_\ell^{-e})$.
Let $\Sigma_{0,\ell}$ consist of infinity and the primes of
$F$ over $3\ell269$.  For every integer $r\ge1$, uniformly
in the integral weight $e$, one has
\[
 \sum_{\ell\mid r}\frac1{b_\ell}
  \log\#H^1(G_{F,\Sigma_{0,\ell}},A_{\ell,e}[r])
 \le C r\bigl(\log D_F+n_F\log(269r)\bigr),             \tag{G44.11}
\]
where $C$ is absolute, $D_F=|\operatorname{disc}F|$ and
$n_F=[F:\mathbb Q]$.

For a scalar strict Selmer group unramified outside
$\Sigma_{0,\ell}$, this is also a valid upper bound for
the normalized trace-kernel correction in (G44.6), even
when the representation $T$ requires a larger ramification
set.  Thus no new estimate for the number of auxiliary
ramified primes of $T$ is needed for this correction.

On the growing host (FG.3), with an actual source rank
$r=k+m-1$, source cutoff $R.t\le t$, and outer cutoff
$t=\lfloor x^{401}\rfloor$, the total correction
for one chosen trace extraction at each prime and each weight
$1\le e<t$ is
\[
 O_\epsilon\!\left(x^{804+41\epsilon/5}(\log x)^2\right)
                 =o\!\left(\exp(x^\epsilon/2)\right).   \tag{G44.12}
\]
This consumes less than any fixed positive fraction of the
allowance in AC-K27 for sufficiently large $x$.
The assertion controls this finite correction; it does not
supply the upper bound for the mixed Selmer group, or an
unbounded multiplicity of separate trace extractions.
\end{corollary}
\begin{proof}
\proofstep{1. A uniform splitting field for the finite coefficient.}
Fix $\ell\mid r$, put $a=v_\ell(r)$ and $q_\ell=\ell^a$.
The finite module $M=A_{\ell,e}[r]$ is killed by $q_\ell$
and is generated by $b_\ell$ elements as an abelian group.
Its order is $q_\ell^{b_\ell}$.
The extension
\[
                 L=F(\sqrt[3]{269},\mu_{q_\ell})       \tag{G44.13}
\]
is Galois over $F$, is unramified outside
$\Sigma_{0,\ell}$, trivializes $M$ for every $e$, and
contains $\mu_{q_\ell}$.  Indeed the first radical kills
the cubic character; over $\mu_{q_\ell}$ the cyclotomic
character is one modulo $q_\ell$, hence its every integral
power is one on $M$.
The Galois assertion uses $\mu_3\subset F$.
For $F_1=F(\sqrt[3]{269})$ we have $[F_1:F]\le3$.
If this degree is three, the integral power basis has
relative discriminant $-27\cdot269^2$, so
\[
 \log D_{F_1}\le3\log D_F+n_F\log(27\cdot269^2).
\]
If the degree is one the same upper bound is immediate.
Tensoring with the cyclotomic order, as in the proof of
(FG.10), and using
$[\mathbb Q(\mu_{q_\ell}):\mathbb Q]\le q_\ell$ and
$\log D_{\mathbb Q(\mu_{q_\ell})}\le q_\ell\log q_\ell$,
gives
\[
 \begin{aligned}
 [L:F]&\le3q_\ell,&n_L&\le3q_\ell n_F,\\
 s_L&\le3q_\ell s_F,&
 \log D_L&\le C_0q_\ell
              \bigl(\log D_F+n_F\log(269q_\ell)\bigr),
 \end{aligned}                                                 \tag{G44.14}
\]
where $s_F=|\Sigma_{0,\ell,\rm fin}|\le3n_F$.
The relative discriminant also proves the stated
ramification support; cyclotomic ramification occurs
only above $\ell$.

Apply AC-K19, (FG.2), with its number of abelian generators
set equal to $b_\ell$.  After division by $b_\ell$, its
right side is
\[
 2\log D_L+2n_L\log8+
 \log q_\ell\,[2([L:F]+n_L+s_L)+n_F+3s_F].
\]
Substituting (G44.14), $s_F\le3n_F$, and
$q_\ell\le r$ bounds this by
$C_1q_\ell(\log D_F+n_F\log(269r))$.
It bounds the logarithm of the global $H^1$ in particular.
The distinct prime powers $q_\ell$ have product $r$ and
sum at most $r$.  For two integers $u,v\ge2$,
$u+v\le uv$; induction proves the assertion for all the
prime-power factors, with equality allowed when there
is only one.  For $r=1$ the sum is empty.
Summing proves (G44.11).

\proofstep{2. The small support suffices for the kernel.}
Let $S(A_{\ell,e})$ be the scalar strict group under
consideration.  Its unramified condition outside
$\Sigma_{0,\ell}$ places it in
$H^1(G_{F,\Sigma_{0,\ell}},A_{\ell,e})$.
Inflation from this group to absolute $H^1$ is injective:
the normal kernel of $G_F\to G_{F,\Sigma_{0,\ell}}$
acts trivially on the coefficient, and the first
inflation--restriction arrow is injective.
If a class is in the kernel of scalar inclusion into
the mixed coefficient, trace shows that its multiple
by $r$ is zero in absolute cohomology.  Injectivity of
inflation shows that the multiple is already zero in
$H^1(G_{F,\Sigma_{0,\ell}},A_{\ell,e})$.
The exact sequence (G44.7), now on this smaller Galois
group, gives the needed lift from
$H^1(G_{F,\Sigma_{0,\ell}},A_{\ell,e}[r])$.
This proves the correction bound without requiring $T$
to factor through the smaller Galois group.
The normalization is exact:
$b_\ell^{-1}\log\#M$ equals the source normalized
$\ell$-length of $M$ multiplied by $\log\ell$.

\proofstep{3. The actual source growth.}
The field degree and discriminant formulas (FG.11),
with the source window inequalities used in (FG.4), give
\[
 n_F=O_\epsilon(x^{2+8\epsilon}),\qquad
 \log D_F=O_\epsilon
       (x^{2+41\epsilon/5}(\log x)^2).
\]
The literal field \texttt{Ranks.rank\_add\_five\_le} in
\texttt{Params.lean} is $r+5\le R.t$, so
$R.t\le\lfloor x^{401}\rfloor$ gives $r\le x^{401}$.
The ranked-host supply audited in AC-G34 even has the
stronger $400$-cutoff reserve; only the displayed weaker
upper bound is used here.
There are fewer than $x^{401}$ weights.
Multiply (G44.11) by this number, use
$\log(269r)=O(\log x)$, and obtain the first estimate
in (G44.12).  For fixed positive $\epsilon$, every fixed
power of $x$ times a fixed power of $\log x$ is
$o(\exp(x^\epsilon/2))$; taking logarithms proves this
last assertion directly.  Source ranks are odd, so the
prime two actually contributes no trace correction,
although the estimate did not need that improvement.
The positive relative slack required to add this error
to a main estimate is exactly the convention in AC-K27.
\end{proof}

\paragraph{The precise source comparison that remains.}
In \path{QuotientHomNonvanishingBound.lean},
\texttt{GoodCongruenceObjects}, lines 473--517, evaluates
\emph{all} block maps at a single selected point $\pi'$.
This is explicitly proved for the center in
\path{CongruenceAnnihilatesCharacterPart.lean},
\texttt{psi\_scal}, lines 191--203:
$\psi_\rho(f\operatorname{Id})=f(\pi')\operatorname{Id}$ and
$\psi_{\chi,c}(f\operatorname{Id})=f(\pi')$ for every $c$.
The two separate transported packets instead give the distinct
maps $\epsilon_0,\epsilon_1$ of (G44.9).
Consequently the common-central-evaluation hypothesis used in
AC-G37-ann cannot be retained merely by replacing one of its
block maps by the other packet's evaluation.
AC-G44-center describes exactly the additional restriction
imposed by requiring a comparison to be balanced over this
paired center.  It does not say that every possible arithmetic
comparison must be balanced in that way.  Restricting the center
to the diagonal scalars restores a common evaluation, but does
not preserve the full Hecke-family congruence comparison by itself.
A new deformation or relative-cohomology construction may retain
more information; none is supplied by (G44.10).

The direction of the existing source theorem is also explicit:
\texttt{QuotientHom}\allowbreak\texttt{NonvanishingBound},
lines 539--549, bounds
\texttt{vcc} above by lengths of \texttt{homRho} and
\texttt{homChi}; \texttt{GlobalLattice}\allowbreak\texttt{IndexTheorem},
\texttt{DecompositionData.lean}, lines 323--336, bounds
\texttt{vcc} above by \texttt{sPi}, the strict Tate lengths,
and stated local and class-group terms.  A small congruence
index therefore does not bound those Selmer lengths.
The remaining usable gate is a theorem controlling the relative
coefficient (G44.8), with compatible integral lattices and local
conditions, in the direction needed for an upper Selmer bound.
If such a theorem is proved, (G44.6) extracts the desired scalar
coefficient with its stated finite-rank correction.  The exact
transport of two separate theta packets alone does not prove it.

\paragraph{Dictionary scope.}
The tensor, dual, and internal-Hom operations here are operations
on representations and their integral lattices.  They are
compatible with the already audited finite-\'etale representation
dictionary in SGA~1, Expos\'e~V, \S\S4 and 7.
No new geometric-square existence theorem, Hodge inequality,
or unverified Grothendieck locator is being asserted.
